\documentclass[12pt]{article}
\usepackage[utf8]{inputenc}
\usepackage[T1]{fontenc}
\usepackage[margin=1in]{geometry}
\usepackage{setspace}
\usepackage{xcolor}
\usepackage[colorlinks=true, urlcolor=blue!60!black, linkcolor=blue!60!black]{hyperref}
\setlength{\parindent}{2em}
\setlength{\parskip}{1.25em}
\renewcommand{\baselinestretch}{1.0}
\begin{document}
\begin{center}
{\LARGE \textit{Gundy v. United States}, 588 U.S. 128 (2019)\par}
\vspace{0.8em}
{\large Case Brief --- Four-Page Edition\par}
\end{center}
\vspace{1em}

\section{Memory Jogger}

An eight-Justice Court upheld SORNA's grant of authority to the Attorney General to specify the Act's applicability to pre-Act sex offenders: Justice Kagan's plurality read the statute to require registration ``as soon as feasible'' and found an intelligible principle, Justice Alito concurred only in the judgment, and three Justices dissented (\href{https://supreme.justia.com/cases/federal/us/588/17-6086/}{Justia}).

\section{Detailed Case Facts}

The Sex Offender Registration and Notification Act of 2006 replaced a patchwork of state registries, under which over 100,000 offenders had escaped registration, with ``a comprehensive national system,'' backed by felony penalties of up to ten years for failing to register after interstate travel. Roughly 500,000 offenders convicted before the Act could not comply with its ordinary timing rule, so \S~20913(d) gave the Attorney General ``the authority to specify the applicability'' of the requirements to them and to prescribe rules for offenders unable to comply. In 2007 the Attorney General applied SORNA in full to all pre-Act offenders. Herman Gundy pleaded guilty in Maryland in 2005 to sexually assaulting a minor; after his release in 2012 he moved to New York without registering.

\section{Procedural History}

Gundy was convicted under 18 U.S.C. \S~2250 in the Southern District of New York after his motion to dismiss on nondelegation grounds failed. The Second Circuit affirmed, joining all eleven circuits to address the issue. The Supreme Court granted certiorari on the nondelegation question; Justice Kavanaugh was confirmed after argument.

\section{Judicial Votes}

Affirmed, \textbf{5--3}, with no majority opinion (decided June 20, 2019).
\begin{itemize}
\item \textbf{Plurality:} Justice Kagan, joined by Justices Ginsburg, Breyer, and Sotomayor.
\item \textbf{Concurring in the judgment:} Justice Alito.
\item \textbf{Dissent:} Justice Gorsuch, joined by Chief Justice Roberts and Justice Thomas.
\item \textbf{Not participating:} Justice Kavanaugh.
\end{itemize}

\section{Holding}

Section 20913(d) does not violate the nondelegation doctrine. The plurality construed it to require the Attorney General to apply SORNA to all pre-Act offenders as soon as feasible---a narrow transitional delegation that easily satisfies the intelligible-principle test. Alito supplied the fifth vote under existing precedent while signaling willingness to reconsider the doctrine with a full Court.

\section{Analysis of Opinions}

\textbf{Plurality (Kagan).} Article I bars delegating powers strictly and exclusively legislative, but Congress may delegate under broad directives so long as it supplies an intelligible principle (\textit{J.W. Hampton}). Gundy read the statute to let the Attorney General require all, some, or none of the pre-Act offenders to register; Kagan conceded that ``[i]f that were so, we would face a nondelegation question.'' But following \textit{Reynolds v. United States}, the declared purpose of a comprehensive system, the past-tense definition of ``sex offender,'' the legislative history, and the provision's title and second clause showed that the discretion concerned only feasibility and timing. So read, the delegation was ``distinctly small-bore,'' far narrower than ``public interest'' standards long upheld; only \textit{Schechter Poultry} and \textit{Panama Refining} (both 1935) had ever failed.

\textbf{Alito (concurring in the judgment).} Since 1935 the Court has upheld extraordinarily capacious standards; he would support reconsidering that approach, but not by singling out this provision.

\textbf{Gorsuch (dissenting).} Drawing on Locke, Montesquieu, bicameralism, and fair notice, Gorsuch would replace the intelligible-principle test with three permissible delegations: filling up details after Congress sets policy, conditional fact-finding, and discretion over matters within another branch's own powers. The intelligible-principle phrase had mutated into a rubber stamp. The plurality invented a feasibility standard the text lacks and united lawmaking and prosecution in the chief prosecutor.

\textbf{Later development.} In \textit{FCC v. Consumers' Research} (2025), the Court, 6--3, reaffirmed the intelligible-principle test; Gorsuch, Thomas, and Alito dissented.

\section{Examples: Future Applications}

\begin{enumerate}
\item \textbf{Same-Side (Upheld): Drone Registry.} An FAA power to specify the registry's applicability to pre-statute operators, read as a feasibility judgment, is small-bore.
\item \textbf{Same-Side (Upheld): Health Standard.} Air standards ``requisite to protect the public health'' were upheld in \textit{Whitman v. American Trucking}.
\item \textbf{Opposite-Side (Likely Invalid): Whether a Crime Applies.} A statute expressly letting the Attorney General decide whether a new firearms felony applies to anyone forecloses a saving construction.
\item \textbf{Opposite-Side (Invalid): Unguided Shipment Ban.} A power to ban any interstate agricultural shipment, with no policy or standard, repeats \textit{Panama Refining}.
\item \textbf{Fence-Sitter: Frontier AI.} A Secretary may create AI felonies as necessary to protect the public interest and national security. The breadth is great, but national-security delegations receive more latitude.
\end{enumerate}

\section{Critique}

\textbf{Avoidance versus fidelity.} As Gorsuch charged, the plurality---and later \textit{Consumers' Research}---saved statutes by narrowing, or effectively rewriting, broad text. Is that respect for Congress or usurpation of its drafting pen?

\textbf{Criminal delegations and structural irony.} \textit{Gundy} applied the same lenient test to a ten-year felony as to utility ratemaking, leaving open the question reserved in \textit{Touby}. And because a workable nondelegation line proved elusive, the Court instead strengthened presidential control over agencies and judicial control over statutory meaning.

\section{Key Quotations}

\begin{quote}
``Given that standard, a nondelegation inquiry always begins (and often almost ends) with statutory interpretation. Only after a court has determined a challenged statute's meaning can it decide whether the law sufficiently guides executive discretion to accord with Article I.''
\end{quote}
\noindent --- Justice Kagan, plurality opinion, 588 U.S. at 135--36.

\begin{quote}
``The Constitution promises that only the people's elected representatives may adopt new federal laws restricting liberty. Yet the statute before us scrambles that design. It purports to endow the nation's chief prosecutor with the power to write his own criminal code governing the lives of a half-million citizens.''
\end{quote}
\noindent --- Justice Gorsuch, dissenting, 588 U.S. at 149--50.

\begin{quote}
``If a majority of this Court were willing to reconsider the approach we have taken for the past 84 years, I would support that effort. But because a majority is not willing to do that, it would be freakish to single out the provision at issue here for special treatment.''
\end{quote}
\noindent --- Justice Alito, concurring in the judgment, 588 U.S. at 149.

\end{document}
